\documentclass[11pt,a4paper]{article}
\usepackage[utf8]{inputenc}
\usepackage[T1]{fontenc}
\usepackage[margin=0.75in]{geometry}
\usepackage{enumitem}
\usepackage{titlesec}
\usepackage{xcolor}
\usepackage{hyperref}

\definecolor{primarycolor}{RGB}{33,37,41}
\definecolor{secondarycolor}{RGB}{0,102,204}
\definecolor{accentcolor}{RGB}{108,117,125}
\hypersetup{colorlinks=true,linkcolor=secondarycolor,urlcolor=secondarycolor}
\titleformat{\section}{\Large\bfseries\color{primarycolor}}{}{0em}{}
\titlespacing*{\section}{0pt}{1.2em}{0.5em}
\setlist[itemize]{leftmargin=*,itemsep=0.2em,topsep=0.25em}
\pagestyle{empty}

\begin{document}

\begin{center}
{\Huge\bfseries\color{primarycolor} Katarzyna Wiśniewska}\\[0.35em]
{\large\color{accentcolor} Pflegehelferin}
\end{center}

\section*{PROFIL}
Pflegehelferin mit fünf Jahren Berufserfahrung als Krankenschwester auf einer Station für Innere Medizin in Polen und praktischer Erfahrung im Seniorenzentrum in Deutschland. Erfahrung mit Menschen mit Demenz, Sturzprophylaxe, der wöchentlichen Leitung einer Singgruppe sowie der Einarbeitung neuer Pflegehelferinnen und Praktikanten.

\section*{BERUFLICHE ERFAHRUNG}
\textbf{Pflegebereich – Praxisprojekte | Seniorenzentrum Am Lindenhof | 2021–heute}
\begin{itemize}
\item Einarbeitung von Pflegehelferinnen und Praktikanten während der ersten zwei Wochen nach dem Einarbeitungsplan der Einrichtung; bisher sieben Personen.
\item Wöchentliche Leitung einer Singgruppe für sechs bis acht Bewohnerinnen und Bewohner mit Demenz; Angehörige lobten die Gruppe im Heimbeirat.
\item Mitarbeit mit der Wohnbereichsleitung an der Sturzprophylaxe: Auswertung von Sturzprotokollen über sechs Monate und Einführung von Nachtlichtern sowie festen Toilettengängen für neun Bewohnerinnen und Bewohner. Die Stürze im Wohnbereich gingen im folgenden Halbjahr von 14 auf 8 zurück.
\end{itemize}

\textbf{Krankenschwester, Station für Innere Medizin | Szpital Opole, Polen | fünf Jahre}
\begin{itemize}
\item Verbandswechsel bei Dekubitus und diabetischem Fuß unter Aufsicht der Stationsärztin (2016–2019).
\end{itemize}

\section*{QUALIFIKATIONEN UND ANERKENNUNG}
\begin{itemize}
\item Derzeitiger beruflicher Status: Pflegehelferin.
\item Anerkennungsverfahren für den polnischen Pflegefachkraftabschluss: Defizitbescheid erhalten; Kenntnisprüfung für März 2027 vorgesehen.
\item LG1-Kurs absolviert.
\item Deutsch: B2-Zertifikat (2021).
\end{itemize}

\section*{KENNTNISSE}
\begin{itemize}
\item Vivendi PD: gute Kenntnisse.
\item MediFox: bisher keine Anwendungserfahrung.
\item Führerschein Klasse B: derzeit nicht vorhanden; Theorieprüfung bestanden, praktische Prüfung im Januar 2027.
\end{itemize}

\section*{ARBEITSZEIT}
\begin{itemize}
\item Teilzeit mit 75 \% bevorzugt.
\item Jedes zweite Wochenende möglich.
\end{itemize}

\end{document}